% Using InitechBase, Introduction (PUB-REF-0001 AN-21 F-ATE: what the
% section is, who should read it, chapter-by-chapter list with chapter
% labels as side heads; annex Sec 8.1).  Folios UI-n.
\SectionIntroduction{Introduction}
\Gates{test-pubs-build}
\doctitle{Using InitechBase} is the reference to InitechBase.  It describes
how commands are put together, how expressions are written, how InitechBase
works with database, index, report form and program files, and the settings
that govern it.  Then it gives every command, every function and every error
message in alphabetical order, with the exact form of each and examples of
its use.

Read \doctitle{Learning InitechBase} first if you have not used InitechBase
before.  Turn to this section when you need the exact form of a command, when
a command does not do what you expected, or when InitechBase displays a
message you do not know.  \doctitle{Programming with InitechBase}, in Volume
Two, shows how the commands are used together in programs; for the form of each
command it refers to this section.

\section{How This Section Is Organized}
\SideHead{Chapter 1}Overview describes databases, records and fields, the
types of field, and the files InitechBase uses.

\SideHead{Chapter 2}Using Commands describes the dot prompt, the parts of a
command, expressions and operators, memory variables, work areas, the record
pointer and index files.

\SideHead{Chapter 3}About Program Files describes how programs are written,
stored and run.

\SideHead{Chapter 4}The InitechBase Environment lists the settings in force
when InitechBase starts and how the SET commands change them.

\SideHead{Chapter 5}Commands gives every command in alphabetical order.

\SideHead{Chapter 6}Functions gives every function in alphabetical order.

\SideHead{Chapter 7}Error Messages lists every message, with its cause and
what to do.

\SideHead{Appendix A}Keys at the Dot Prompt.

\SideHead{Appendix B}Beginning and End of File.

\SideHead{Appendix C}Structure of InitechBase Files.

\SideHead{Appendix D}ASCII Character Codes.

\SideHead{Appendix E}Data Types in Expressions.

\SideHead{Glossary}The terms used in this manual.

\section{The Examples}
Every example in this section was run against a fresh copy of the files on the
Sample Programs and Utilities diskette, described in Chapter~1 of
\doctitle{Getting Started}.  When you try an example yourself, work on a copy
of the diskette, because some examples change the files.  A record number,
total or date in your results is the same as in the example only if your files
are as they were on the diskette.

The examples were run on the 1st of March 1992; \fn{DATE()} returns that
date in them.  Your computer's date takes its place when you try them.
